\section{Introduction}
\label{sec:introduction}

Embodied policy learning increasingly depends on the quality and downstream
utility of demonstrations, rather than on record count alone
\citep{yeDataPyramidEmbodied2026,hejnaRobotDataCuration2025,hoqueEGODEXLEARNINGDEXTEROUSMANIPULATION2026,collaborationOpenXEmbodimentRobotic2025}.
Simulation, portable teleoperation, egocentric recording, and large-scale
real-robot datasets have broadened the sources of manipulation data
\citep{collaborationOpenXEmbodimentRobotic2025,chiUniversalManipulationInterface2024,khazatskyDROIDLargeScaleInTheWild2024,hoqueEGODEXLEARNINGDEXTEROUSMANIPULATION2026,mandlekarMimicGenDataGeneration2023}.
Yet these collection paradigms do not share a common record of temporal
synchronization, task phase, operator correction, and failure cause. Recent
visual auditing and reward-modeling methods begin to represent failures,
progress, near misses, and partial completion explicitly
\citep{duanAHAVisionLanguageModelDetecting2024,leeRoboRewardGeneralPurposeVisionLanguage2026,sermanetRoboVQAMultimodalLongHorizon2024,wuLargeRewardModels2026}.
For visually observable manipulation with tight action tolerances,
misalignment, stagnation, and repeated attempts reduce the reuse value of a
demonstration. Data curation and visuomotor policy studies likewise show that
trajectory properties and data selection affect downstream learning
\citep{hejnaRobotDataCuration2025,chiDiffusionPolicyVisuomotor2024,zhaoLearningFineGrainedBimanual2023}.

High-quality demonstrations do not require assistance at every step. Smooth,
single-pass executions preserve the operator's nominal behavior, while local
corrections reveal how an expert recovers from difficult states
\citep{rossReductionImitationLearning2011,liuRobotLearningJob2023,luoPreciseDexterousRobotic}.
Current work covers the required components separately: teleoperation systems
address how demonstrations are collected
\citep{wuGELLOGeneralLowCost2024,chengOpenTeleVisionTeleoperationImmersive2024,wangDexCapScalablePortable2024,qinAnyTeleopGeneralVisionBased2024,xuDexUMIUsingHuman2025};
data-quality and visual models score trajectories after collection
\citep{hejnaRobotDataCuration2025,duanAHAVisionLanguageModelDetecting2024,leeRoboRewardGeneralPurposeVisionLanguage2026};
and shared control, residual learning, and interactive imitation learning
primarily optimize deployment success or expert-intervention cost
\citep{johanninkResidualReinforcementLearning2019,javdaniSharedAutonomyHindsight2015,hoqueThriftyDAggerBudgetAwareNovelty2021,welteFlowCorrectEfficientInteractive2026a}.
What remains unclear is whether audited historical demonstrations can improve
the yield of the \emph{next} teleoperation round under fixed operator-time and
safety budgets, while preserving operator commands in nominal states.

We study this question through a causal, visual-conditioned continuous action
filter. Correction intervals provide weak supervision for difficult portions
of a trajectory; successful nominal segments constrain unnecessary action
offsets. A frozen visual encoder supplies continuous observations, and an
offline vision-language auditor proposes phase, progress, stagnation, and
recovery events for human review. The online model predicts a candidate expert
action and a time-varying assistance gain. Their difference from the raw human
command produces a bounded residual whose strength changes gradually with task
context. Executing the filtered command produces the next visual observation,
and newly audited episodes supply the next training round. Fig.~\ref{fig:overview}
summarizes this loop.

\begin{figure*}[t]
  \centering
  \resizebox{0.99\textwidth}{!}{%
  \begin{tikzpicture}[
    font=\sffamily\footnotesize,
    >=Latex,
    flow/.style={->,line width=0.8pt,draw=DarkGray},
    data/.style={draw=AuditBlue,fill=AuditBlue!7,rounded corners=2pt,
      minimum height=10mm,align=center,inner xsep=5pt},
    model/.style={draw=ActionTeal,fill=ActionTeal!7,rounded corners=2pt,
      minimum height=10mm,align=center,inner xsep=5pt},
    gain/.style={draw=GainOrange,fill=GainOrange!9,rounded corners=2pt,
      minimum height=10mm,align=center,inner xsep=5pt},
    safety/.style={draw=SafetyRed,fill=SafetyRed!7,rounded corners=2pt,
      minimum height=10mm,align=center,inner xsep=5pt},
    tag/.style={font=\sffamily\scriptsize\bfseries,text=DarkGray}
  ]
    % Offline audit and update path.
    \node[data,minimum width=28mm] (episodes) at (0,2.55)
      {Synchronized\\episodes};
    \node[data,minimum width=31mm,right=7mm of episodes] (audit)
      {Offline VLM proposals\\+ human review};
    \node[data,minimum width=33mm,right=7mm of audit] (views)
      {Training views\\nominal + correction};
    \node[model,minimum width=29mm,right=7mm of views] (train)
      {Filter training\\and selection};
    \draw[flow] (episodes) -- (audit);
    \draw[flow] (audit) -- (views);
    \draw[flow] (views) -- (train);
    \node[tag,anchor=west] at ($(episodes.north west)+(0,4mm)$)
      {OFFLINE AUDIT AND UPDATE};

    % Online input path.
    \node[data,minimum width=24mm] (human) at (0,0)
      {Operator command\\$u_t^{\mathrm{raw}}$};
    \node[data,minimum width=24mm,below=5mm of human] (vision)
      {Multi-view image\\$o_t$};
    \node[data,minimum width=24mm,below=5mm of vision] (state)
      {Robot state\\$q_t$};

    \node[model,minimum width=28mm] (context) at (4.25,-0.75)
      {Frozen visual encoder\\+ causal history};
    \node[model,minimum width=28mm,right=7mm of context] (backbone)
      {Causal temporal\\backbone $h_t$};
    \node[model,minimum width=30mm,right=7mm of backbone,yshift=8mm] (action)
      {CVAE action branch\\$\hat u_t^{\mathrm{exp}}$};
    \node[gain,minimum width=30mm,right=7mm of backbone,yshift=-8mm] (gain)
      {Rate-limited gain\\$\alpha_t$};
    \node[model,minimum width=29mm,right=8mm of action,yshift=-8mm] (compose)
      {Bounded residual\\$u_t^{\mathrm{raw}}+\alpha_t\hat\delta_t$};
    \node[safety,minimum width=24mm,right=7mm of compose] (safe)
      {Safety\\projection};
    \node[data,minimum width=24mm,right=7mm of safe] (robot)
      {Robot\\execution};

    \draw[flow] (human.east) -- (context.west);
    \draw[flow] (vision.east) -- (context.west);
    \draw[flow] (state.east) -- (context.west);
    \draw[flow] (context) -- (backbone);
    \draw[flow] (backbone.east) -- ++(4mm,0) |- (action.west);
    \draw[flow] (backbone.east) -- ++(4mm,0) |- (gain.west);
    \draw[flow] (action.east) -- ++(4mm,0) |- (compose.west);
    \draw[flow] (gain.east) -- ++(4mm,0) |- (compose.west);
    \draw[flow] (compose) -- (safe);
    \draw[flow] (safe) -- node[above,tag]{$u_t^{\mathrm{out}}$} (robot);

    \draw[flow,draw=AuditBlue]
      (robot.south) -- ++(0,-9mm) -| node[pos=0.25,below,tag]
      {next observation $o_{t+1}$} (vision.south);
    \draw[flow,draw=AuditBlue,dashed]
      (robot.north) -- ++(0,24mm) -| (episodes.north);
    \node[tag,anchor=south] at ($(audit.north)!0.5!(views.north)+(0,8mm)$)
      {new candidate episodes $D'_r$};
    \draw[flow,draw=ActionTeal,dashed]
      (train.south) -- ++(0,-7mm) -| (backbone.north);

    \node[tag,anchor=west] at ($(human.north west)+(0,5mm)$)
      {ONLINE CONTINUOUS ASSISTANCE};
    \begin{scope}[on background layer]
      \node[draw=DarkGray!35,fill=SoftGray!45,rounded corners=3pt,
        fit=(episodes)(audit)(views)(train),inner sep=4mm] {};
      \node[draw=DarkGray!35,fill=SoftGray!25,rounded corners=3pt,
        fit=(human)(vision)(state)(context)(backbone)(action)(gain)(compose)(safe)(robot),
        inner sep=4mm] {};
    \end{scope}
  \end{tikzpicture}}
  \caption{System overview. Offline auditing converts synchronized episodes
  into complementary nominal and correction supervision. During collection,
  the causal filter predicts a candidate expert action and a continuous
  assistance gain; bounded residual composition and safety projection produce
  the executed command. New executions are audited before they update the next
  collection round. Dashed arrows indicate cross-round information flow.}
  \label{fig:overview}
\end{figure*}


Our contributions are threefold:
\begin{enumerate}
  \item We introduce a correction-aware, auditable demonstration
  representation in which correction intervals and successful nominal segments
  provide complementary supervision, while separate action and audit sets make
  training admission explicit.
  \item We propose a visual-conditioned continuous action filter that models a
  candidate expert action separately from assistance gain, and constrains both
  the magnitude and rate of assistance.
  \item We formulate and evaluate an iterative real-robot collection loop whose
  primary outcome is valid-demonstration yield under a fixed operator budget,
  with independent auditing and downstream imitation learning used to assess
  reliability and utility.
\end{enumerate}
